% ═══════════════════════════════════════════════════════════════
% MT-02c-ML: Musterlösung Minitest Possessivbegleiter
% Spanisch Klasse 7
% ═══════════════════════════════════════════════════════════════
% Dauer: 10 Minuten
% Punkte: 26
% Inhalt: mi/tu/su, Singular zu Plural, Dialog
% ═══════════════════════════════════════════════════════════════

\documentclass[11pt, a4paper]{article}

\usepackage[utf8]{inputenc}
\usepackage[T1]{fontenc}
\usepackage[ngerman]{babel}
\usepackage{helvet}
\renewcommand{\familydefault}{\sfdefault}

\usepackage[margin=2cm]{geometry}
\usepackage{enumitem}
\usepackage{tabularx}
\usepackage{booktabs}

\usepackage{xcolor}
\usepackage{tcolorbox}
\tcbuselibrary{skins}

\usepackage{fancyhdr}
\usepackage{lastpage}
\usepackage{needspace}

% FARBEN
\definecolor{spanisch-color}{HTML}{c41e3a}
\definecolor{grau-color}{HTML}{888888}
\definecolor{hellgrau-color}{HTML}{f5f5f5}
\definecolor{loesung-color}{HTML}{c62828}

\colorlet{spanisch}{spanisch-color}
\colorlet{grau}{grau-color}
\colorlet{hellgrau}{hellgrau-color}
\colorlet{loesung}{loesung-color}

\newcommand{\fachfarbe}{spanisch}

\newcommand{\thema}{Minitest: Possessivbegleiter -- MUSTERLÖSUNG}
\newcommand{\fach}{Spanisch}
\newcommand{\klasse}{7}
\newcommand{\maxpunkte}{26}
\newcommand{\zeit}{10 Minuten}
\newcommand{\datum}{\today}

\pagestyle{fancy}
\fancyhf{}
\renewcommand{\headrulewidth}{0.4pt}
\renewcommand{\footrulewidth}{0.4pt}

\lhead{\fach{} -- Klasse \klasse}
\chead{\textbf{MT-02c -- ML}}
\rhead{\datum}

\lfoot{\zeit}
\cfoot{}
\rfoot{Seite \thepage\ von \pageref{LastPage}}

\newcommand{\punktebox}[1]{\hfill\fbox{\small \_/#1}}
\newcommand{\luecke}[1]{\rule{#1}{0.4pt}}
\newcommand{\lsg}[1]{\textcolor{loesung}{\textbf{#1}}}

\newtcolorbox{infobox}{
    colback=hellgrau,
    colframe=\fachfarbe,
    boxrule=1pt,
    arc=0pt,
    left=8pt, right=8pt, top=8pt, bottom=8pt
}

\newenvironment{aufgabe}[1]{%
    \needspace{5\baselineskip}%
    \nopagebreak[4]%
    \vspace{10pt}
    \noindent\textbf{\textcolor{\fachfarbe}{Aufgabe #1}}
    \nopagebreak[4]%
    \vspace{5pt}
    \nopagebreak[4]%
    \begin{list}{}{\leftmargin=0pt}
    \item[]
}{%
    \end{list}
}

\newcommand{\es}[1]{\textcolor{\fachfarbe}{\textbf{#1}}}

\begin{document}

% ═══════════════════════════════════════════════════════════════
% KOPFBEREICH
% ═══════════════════════════════════════════════════════════════

\begin{center}
    {\Large\bfseries\textcolor{\fachfarbe}{\thema}}
\end{center}

\vspace{5pt}

\noindent
\begin{tabularx}{\textwidth}{|X|X|X|}
\hline
\textbf{Name:} \lsg{LÖSUNG} &
\textbf{Klasse:} \klasse &
\textbf{Datum:} \luecke{2cm} \\
\hline
\end{tabularx}

\vspace{10pt}

\begin{infobox}
\textbf{Zeit:} \zeit{} | \textbf{Punkte:} \maxpunkte{} | \textbf{Hilfsmittel:} keine
\end{infobox}

\vspace{15pt}

% ═══════════════════════════════════════════════════════════════
% AUFGABE 1: mi, tu, su (8 Punkte)
% ═══════════════════════════════════════════════════════════════

\begin{aufgabe}{1 -- mi, tu oder su?}
Ergänze den passenden Possessivbegleiter: \es{mi}, \es{tu} oder \es{su}. \punktebox{8}
\end{aufgabe}

\noindent
1. \lsg{Mi} madre se llama Carmen. \hfill (meine Mutter)

\vspace{0.8em}

\noindent
2. \lsg{Tu} padre es profesor. \hfill (dein Vater)

\vspace{0.8em}

\noindent
3. \lsg{Su} hermano tiene 15 años. \hfill (sein Bruder)

\vspace{0.8em}

\noindent
4. \lsg{Tu} abuela es muy simpática. \hfill (deine Großmutter)

\vspace{0.8em}

\noindent
5. \lsg{Mi} tío vive en Madrid. \hfill (mein Onkel)

\vspace{0.8em}

\noindent
6. \lsg{Su} prima se llama Ana. \hfill (ihre Cousine)

\vspace{0.8em}

\noindent
7. \lsg{Mi} hermana es muy alta. \hfill (meine Schwester)

\vspace{0.8em}

\noindent
8. \lsg{Tu} abuelo tiene 70 años. \hfill (dein Großvater)

\vspace{15pt}

% ═══════════════════════════════════════════════════════════════
% AUFGABE 2: Singular zu Plural (8 Punkte)
% ═══════════════════════════════════════════════════════════════

\begin{aufgabe}{2 -- Singular zu Plural}
Bilde die Pluralform. Denke daran: mi $\rightarrow$ mis, tu $\rightarrow$ tus, su $\rightarrow$ sus! \punktebox{8}
\end{aufgabe}

\noindent
a) \es{mi hermano} $\rightarrow$ \lsg{mis hermanos}

\vspace{0.8em}

\noindent
b) \es{tu abuelo} $\rightarrow$ \lsg{tus abuelos}

\vspace{0.8em}

\noindent
c) \es{su prima} $\rightarrow$ \lsg{sus primas}

\vspace{0.8em}

\noindent
d) \es{mi tío} $\rightarrow$ \lsg{mis tíos}

\vspace{15pt}

% ═══════════════════════════════════════════════════════════════
% AUFGABE 3: Dialog (10 Punkte)
% ═══════════════════════════════════════════════════════════════

\begin{aufgabe}{3 -- Dialog vervollständigen}
Ergänze den Dialog mit dem passenden Possessivbegleiter. \punktebox{10}
\end{aufgabe}

\noindent
\textbf{A:} ¿Cómo se llama \lsg{tu} madre?

\vspace{0.8em}

\noindent
\textbf{B:} \lsg{Mi} madre se llama Carmen.

\vspace{0.8em}

\noindent
\textbf{A:} ¿Y \lsg{tus} hermanos? ¿Cómo se llaman?

\vspace{0.8em}

\noindent
\textbf{B:} \lsg{Mis} hermanos se llaman Pablo y Ana.

\vspace{0.8em}

\noindent
\textbf{A:} ¿Cuántos años tiene \lsg{tu} padre?

% ═══════════════════════════════════════════════════════════════
% PUNKTEÜBERSICHT
% ═══════════════════════════════════════════════════════════════

\vspace{25pt}
\hrule
\vspace{10pt}

\noindent
\begin{tabularx}{\textwidth}{|X|c|c|c||c|}
\hline
\textbf{Aufgabe} & 1 & 2 & 3 & \textbf{Gesamt} \\
\hline
\textbf{max. Punkte} & 8 & 8 & 10 & \textbf{26} \\
\hline
\textbf{erreicht} & \lsg{8} & \lsg{8} & \lsg{10} & \lsg{26} \\[0.8em]
\hline
\end{tabularx}

\vspace{10pt}

\begin{flushright}
\textbf{Note:} \lsg{14 NP}
\end{flushright}

\end{document}
